\documentclass[a4paper,12pt]{article}
\usepackage[utf8]{inputenc}
\usepackage[T1]{fontenc}
\usepackage[english]{babel}
\usepackage{geometry}
\geometry{left=2.5cm,right=2.5cm,top=2.5cm,bottom=2.5cm}

\usepackage{lmodern}
\usepackage{xcolor}
\definecolor{titleblue}{RGB}{0,51,140}

\usepackage{hyperref}

\pagestyle{empty}

\begin{document}

\begin{flushleft}
    \textbf{Ismail AISSAOUI} \\
    State Engineer in Artificial Intelligence \\
    ismail.aissaoui.pro@gmail.com \\
    +213 660 70 77 96 \\
    \href{https://ismail-aissaoui.vercel.app}{ismail-aissaoui.vercel.app}
\end{flushleft}

\begin{flushright}
    To the PhD Selection Committee \\
    \textbf{EDT Program - Inria Rennes} \\
    Attn: Prof. Benoit Caillaud \\
    \medskip
    May 1, 2026
\end{flushright}

\bigskip

\noindent \textbf{Subject: Application for PhD: Structural Methods for Mixed Model/Data Digital Twin Engineering (Rank 2 - PC1)}

\bigskip

Dear Prof. Caillaud,

As a State Engineer in Artificial Intelligence from ESI-SBA, I am writing to express my enthusiastic application for the PhD position in "Structural Methods for Mixed Model/Data Digital Twin Engineering" within the Engineering Digital Twin (EDT) national program. This application is motivated by the ideal continuation this position represents for my work in hybrid physics-data modeling.

Currently, I am completing my thesis at Sonatrach, where I have developed a **Physics-Informed Digital Twin** for gas turbine monitoring. My research focuses on the integration of physical laws (thermodynamics) into deep learning architectures to bridge the gap between pure simulation and real-world sensor data. This experience has given me a deep appreciation for the challenges of **model/data mismatches** and the need for robust diagnostic tools to ensure the reliability of digital twins.

I am particularly interested in your proposal to use **structural analysis of differential-algebraic equation (DAE) systems** to detect inconsistencies in large-scale models. My background in **Graph Neural Networks (GNNs)** and relational modeling provides me with a solid foundation for the graph-based algorithms required to compute Minimal Structurally Overdetermined (MSO) subsystems. I am eager to apply these mathematical techniques to localize deficiencies in equation-based models and assist in targeted model correction.

My background in graph algorithms and my experience with production-grade PyTorch and C++ implementations provide me with the necessary tools to contribute to the **Artemis platform** and the broader EDT ecosystem. I am a highly autonomous researcher, comfortable bridging the gap between mathematical structural analysis and industrial physical simulations.

Thank you for your time and consideration. I look forward to the possibility of discussing how my experience in hybrid modeling and relational data structures can contribute to the success of your research group at Inria.

Sincerely,

\bigskip

\textbf{Ismail AISSAOUI}

\end{document}
